% [原文 p.67]
\[
\operatorname{totaldiscrep}(X,\Delta):=\inf\{a(E,X,\Delta)\},
\]
其中 $E$ 取遍 $X$ 上空的一切除子（即取遍一切双有理态射 $f\colon Y\to X$ 的一切例外素除子，以及 $X$ 的一切素除子）。等价地说，$\operatorname{totaldiscrep}(X,\Delta):=\inf\{a(v,X,\Delta)\}$，其中 $v$ 取遍 $X$ 上的一切代数赋值。
通常不用 $\operatorname{discrep}(X,0)$，而简写为 $\operatorname{discrep}(X)$；$\operatorname{totaldiscrep}(X)$ 也类似。\hfill$\square$

利用下面两个引理，在多数情形下都可以计算差异系数。证明都容易。第一个引理说明一次爆破时差异系数如何变化，第二个引理说明如何归纳地使用它。$X$ 上空的任意除子都可用此法处理，这一事实的依据是引理 2.45。

\begin{lemma}{2.29}
设 $X$ 为非奇异簇，$\Delta=\sum a_iD_i$ 为互不相同素除子之和。又设 $Z\subset X$ 为余维数 $k$ 的闭子簇。设 $p\colon B_ZX\to X$ 为以 $Z$ 为中心的 $X$ 的爆破；在例外除子的不可约分支中，$E\subset B_ZX$ 是支配 $Z$ 的那一个。（若 $Z$ 非奇异，则除 $E$ 之外没有例外除子。）此时成立
\[
a(E,X,\Delta)=k-1-\sum_i a_i\operatorname{mult}_ZD_i,
\]
其中 $\operatorname{mult}_ZD_i$ 是 $D_i$ 在 $Z$ 的泛点处的重数。
\end{lemma}

\begin{lemma}{2.30}
设 $f\colon Y\to X$ 为正规簇之间的固有双有理态射。$\Delta_Y$（或 $\Delta_X$）为 $Y$（或 $X$）上的 $\Q$-除子，满足
\[
K_Y+\Delta_Y\equiv f^*(K_X+\Delta_X),\qquad f_*\Delta_Y=\Delta_X.
\]
此时对 $X$ 上空任意除子 $F$，成立
\[
a(F,Y,\Delta_Y)=a(F,X,\Delta_X).
\]
\end{lemma}

下例是能计算差异系数的第一个例子。

\begin{corollaryp}{2.31}
设 $\Delta=\sum a_iD_i$ 为 $\Q$-除子。此时
\begin{enumerate}[label=（\arabic*）,leftmargin=2.6em,itemsep=1pt]
\item $\operatorname{discrep}(X,\Delta)=-\infty$，或
\[
-1\le\operatorname{totaldiscrep}(X,\Delta)\le\operatorname{discrep}(X,\Delta)\le1;
\]
% [原文 p.68]
\item 若 $X$ 非奇异，则 $\operatorname{discrep}(X,0)=1$；
\item 设 $X$ 非奇异，$\sum D_i$ 为单纯正规交叉除子，且对一切 $i$ 有 $a_i\le1$。此时
\[
\operatorname{discrep}(X,\Delta)=\min\Bigl\{\min_{i\ne j,\ D_i\cap D_j\ne\varnothing}\{1-a_i-a_j\},\ \min_i\{1-a_i\},\ 1\Bigr\}.
\]
\end{enumerate}
\end{corollaryp}
\begin{proof}
对既不含于 $\Supp X$、也不含于 $X$ 的奇点集的余维数 $2$ 的子簇作爆破，可知
\[
\operatorname{discrep}(X,\Delta)\le1.
\]
设存在 $X$ 上空的除子 $E$ 使 $a(E,X,\Delta)<-1$。取 $c>0$ 使 $a(E,X,\Delta)=-1-c$。取双有理态射 $f\colon Y\to X$，使 $E$ 成为正规簇 $Y$ 上的除子，并按记号 2.26 写出公式 $K_Y+\Delta_Y\equiv f^*(K_X+\Delta)$。
取 $Y$ 的余维数 $2$ 的闭子簇 $Z_0$，使它含于 $E$，但不含于 $f$ 的其它例外除子、$\Supp f_*^{-1}\Delta$ 以及 $Y$ 的奇点集。以 $Z_0$ 为中心的爆破记为 $g_1\colon Y_1=B_{Z_0}Y\to Y$，$E_1\subset Y_1$ 为支配 $Z_0$ 的例外除子。此时
\[
a(E_1,X,\Delta)=a(E_1,Y,\Delta_Y)=-c.
\]
设 $Z_1\subset Y_1$ 为 $E$ 的双有理变换与 $E_1$ 的公共部分。以 $Z_1$ 为中心的爆破记为 $g_2\colon Y_2=B_{Z_1}Y_1\to Y_1$，支配 $Z_1$ 的例外除子为 $E_2\subset Y_2$。此时
\[
a(E_2,X,\Delta)=a(E_2,Y,\Delta_Y)=-2c.
\]
这次沿 $E$ 的双有理变换与 $E_2$ 的公共部分作爆破，得到差异系数为 $-3c$ 的除子。反复进行，便得到 $-4c,-5c,\dots$。由此得到 (1)。
图 2.2 是上述论证的说明。
\begin{figure}[htbp]
\centering
\resizebox{0.72\textwidth}{!}{\begin{tikzpicture}[line width=0.5pt,font=\footnotesize]
% 第一栏
\draw (0.85,3.42) -- (0.30,2.58);
\draw (0.80,2.62) -- (0.30,1.82);
\draw (0.80,1.86) -- (0.30,1.06);
\draw (0.85,1.10) -- (0.30,0.30);
\node[anchor=west] at (1.00,3.02) {$E\ (-1-c)$};
\node[anchor=west] at (1.00,2.27) {$E_3\ (-3c)$};
\node[anchor=west] at (1.00,1.52) {$E_2\ (-2c)$};
\node[anchor=west] at (1.00,0.77) {$E_1\ (-c)$};
% 第二栏
\draw (3.65,3.42) -- (3.10,2.58);
\draw (3.60,2.62) -- (3.10,1.82);
\draw (3.60,1.86) -- (3.10,1.06);
\node[anchor=west] at (3.80,3.02) {$E\ (-1-c)$};
\node[anchor=west] at (3.80,2.27) {$E_2\ (-2c)$};
\node[anchor=west] at (3.80,1.52) {$E_1\ (-c)$};
% 第三栏
\draw (6.45,3.42) -- (5.90,2.58);
\draw (6.40,2.62) -- (5.90,1.82);
\node[anchor=west] at (6.60,3.02) {$E\ (-1-c)$};
\node[anchor=west] at (6.60,2.27) {$E_1\ (-c)$};
% 第四栏
\draw (9.25,3.42) -- (8.70,0.30);
\node[anchor=west] at (9.40,3.02) {$E\ (-1-c)$};
\fill (9.02,1.90) circle (0.05);
\node[anchor=west] at (9.15,1.80) {$Z_0$ 的泛点};
% 收缩箭头
\draw[->] (2.55,2.45) -- (2.95,2.45);
\node[above] at (2.75,2.45) {$g_3$};
\draw[->] (5.35,2.45) -- (5.75,2.45);
\node[above] at (5.55,2.45) {$g_2$};
\draw[->] (8.15,2.45) -- (8.55,2.45);
\node[above] at (8.35,2.45) {$g_1$};
\end{tikzpicture}
}
\figtitle{2.2}{差异系数的值（原书 p.69）}
\end{figure}
(2) 是 (3) 的特殊情形。
设 $r(X,\Delta)$ 为 (3) 中待证等式的右端。爆破 $D_i\cap D_j$，可知 $\operatorname{discrep}(X,\Delta)\le r(X,\Delta)$。设 $D$ 为某双有理态射 $f\colon Y\to X$ 的例外除子。须证 $a(D,X,\Delta)\ge r(X,\Delta)$。由引理 2.45，$D$ 可由以非奇异闭子簇为中心的有限次爆破得到
% [原文 p.69]
（$t$ 次，只需每次限制在稠密开集上即可）。用关于 $t$ 的归纳法证明不等式。注意 $r(X,\Delta)$ 限制到稠密开集上不减少。因此不妨设 $f(D)$ 是 $X$ 的非奇异闭子簇。设 $g_1\colon X_1\to X$ 为以 $f(D)$ 为中心的爆破，$E_1\subset X_1$ 为例外除子。把 $X$ 限制到含 $f(D)$ 的一般点的开集上，可设 $E_1\cup(g_1)_*^{-1}\Delta$ 为单纯正规交叉除子。同样，把 $Y$ 限制到含 $D$ 的一般点的开集上，可设 $f_1\colon Y\to X_1$ 为态射。适当取整数 $b\le k$ 并重新编号 $D_i$（$i\ge1$），可使 $\operatorname{codim}f(D)=k\ge2$，且 $f(D)\subset D_i$ 恰好对 $i\le b$ 成立。由爆破公式，
\[
a(E_1,X,\Delta)=k-1-\sum_{i\le b}a_i.
\]
分 $b\le0$、$b=1$、$b\ge2$ 三种情形讨论。
若 $b\le0$，则 $a(E_1,X,\Delta)\ge1\ge r(X,\Delta)$。
若 $b=1$，则 $a(E_1,X,\Delta)\ge1-a_1\ge r(X,\Delta)$。
若 $b\ge2$，则
\[
a(E_1,X,\Delta)\ge(k-b-1)+\sum_{1\le i\le b}(1-a_i)\ge-1+(1-a_1)+(1-a_2)\ge r(X,\Delta).
\]
这样 $t=1$ 的情形得证。另一方面，若用 $K_{X_1}+\Delta_1\equiv g_1^*(K_X+\Delta)$ 定义 $X_1$ 上的 $\Delta_1$，则
% [原文 p.70]
\[
r(X_1,\Delta_1)\ge\min\Bigl\{r(X,\Delta),\ 1+a(E_1,X,\Delta)-\max_{D_i\cap f(D)\ne\varnothing}a_i\Bigr\}
\]
\[
\ge\min\{r(X,\Delta),\ a(E_1,X,\Delta)\}\ge r(X,\Delta).
\]
由于 $\Supp\Delta_1$ 为正规交叉除子，且 $\Delta_1$ 中 $E_1$ 的系数 $-a(E_1,X,\Delta)\le1$，故对 $f_1$ 用归纳假设，即可证明 $a(D,X,\Delta)\ge r(X_1,\Delta_1)\ge r(X,\Delta)$。
\end{proof}

% [Tutor Note: 差异系数总有上界 1，这正是把非奇异簇的值规范为 1 的原因。凡出现 $>1$ 的情形都对应极特殊的奇点，而 $<-1$ 一直下探则对应「无穷坏」的奇点。]

举出可由奇点消解得到的例外除子来计算差异系数的情形。

\begin{corollaryp}{2.32}
设 $f\colon Y\to X$ 为奇点消解，$E\subset Y$ 为例外集，$E_i\subset E$ 取遍一切例外素除子。
\begin{enumerate}[label=（\arabic*）,leftmargin=2.6em,itemsep=1pt]
\item 假设 $\min_i\{a(E_i,X)\}\ge0$。此时成立
\[
\operatorname{discrep}(X)=\min\bigl\{\min_i\{a(E_i,X)\},\ 1\bigr\}.
\]
\item 设 $\Delta=\sum a_iD_i$ 为互不相同素除子之和的表示，且系数满足 $a_i\le1$。此时可以取到 $(X,\Delta)$ 的对数奇点消解 $f$，使 $\sum f_*^{-1}D_i$ 非奇异。任取这样的 $f$。若 $\min_i\{a(E_i,X,\Delta)\}\ge-1$，则成立
\[
\operatorname{discrep}(X,\Delta)=\min\bigl\{\min_i\{a(E_i,X,\Delta)\},\ \min_j\{1-a_j\},\ 1\bigr\}.
\]
\end{enumerate}
\end{corollaryp}
\begin{proof}
关于 (1)：取使 $K_Y+\Delta_Y\equiv f^*K_X$ 成立的 $Y$ 上的 $\Q$-除子 $\Delta_Y$。此时 $\Delta_Y=-\sum_i a(E_i,X)E_i\le0$，且由推论 2.31，$\operatorname{discrep}(Y,\Delta_Y)\ge\operatorname{discrep}(Y,0)=1$；特别地，$\operatorname{discrep}(Y,\Delta_Y)=1$。因此
\[
\operatorname{discrep}(X)=\min\bigl\{\operatorname{discrep}(Y,\Delta_Y),\ \min_i\{a(E_i,X)\}\bigr\}=\min\bigl\{\min_i\{a(E_i,X)\},\ 1\bigr\}.
\]
为说明 (2) 的存在性，取 $(X,\Delta)$ 的任意对数奇点消解 $g\colon Z\to X$。此时 $\operatorname{Sing}\bigl(\sum g_*^{-1}D_i\bigr)$ 是余维数 $2$ 的非奇异子簇
% [原文 p.71]
（设为 $k(g)$ 个）之并。设 $h\colon Z'\to Z\to X$ 为沿各不可约分支 $S$ 的爆破，则 $k(h)=k(g)-1$。因此对 $k(g)$ 归纳可得 $k(f)=0$，即得到使 $\sum f_*^{-1}D_i$ 非奇异的 $f$。设 $\Delta_Y$ 为满足 $K_Y+\Delta_Y\equiv f^*(K_X+\Delta)$ 的 $Y$ 上的 $\Q$-除子，令 $b_i=-a(E_i,X,\Delta)$。由推论 2.31，$\operatorname{discrep}(Y,\Delta_Y)$ 是下列形状之值的最小值：$1-b_i$、$1-b_i-b_j$、$1-a_j$、$1-b_i-a_j$、$1$。（因 $f_*^{-1}D_i$ 与 $f_*^{-1}D_j$ 不相交，故不必考虑 $1-a_j-a_{j'}$ 形状的数。）进一步，
\[
\operatorname{discrep}(X,\Delta)=\min\bigl\{\operatorname{discrep}(Y,\Delta_Y),\ \min_i\{a(E_i,X,\Delta)\}\bigr\}=\min\bigl\{\operatorname{discrep}(Y,\Delta_Y),\ \min_i\{-b_i\}\bigr\}.
\]
这里注意 $-b_i\le1-b_i-b_j$ 与 $-b_i\le1-b_i-a_j$，于是
\[
\operatorname{discrep}(X,\Delta)=\min\bigl\{\min_j\{1-a_j\},\ \min_i\{-b_i\},\ 1\bigr\}.
\]
\end{proof}

下一个结果表明差异系数是下半连续函数，这是最初的一例。这类关于差异系数的比较将在 §5.4 正式进行。

\begin{corollaryp}{2.33}
取对 $(X,\Delta)$ 与 $X$ 上的线性系 $|L|$。设 $L_0\in|L|$ 为任意元，$L_g\in|L|$ 为一般元。此时
\[
\operatorname{discrep}(X,\Delta+cL_0)\le\operatorname{discrep}(X,\Delta+cL_g).
\]
\end{corollaryp}
\begin{proof}
取具有下列性质的奇点消解 $f\colon Y\to X$。
\begin{enumerate}[label=（\arabic*）,leftmargin=2.6em,itemsep=1pt]
\item $f^*|L|=B+|F|$，其中 $|F|$ 为无基点线性系；
\item $f_*^{-1}\Supp\Delta$ 非奇异（参见推论 2.32）；
\item $B+f_*^{-1}\Supp\Delta+\Exc(f)$ 为单纯正规交叉除子。
\end{enumerate}
对任意 $L_\lambda\in|L|$，考虑公式
\[
f^*(K_X+\Delta+cL_\lambda)\equiv K_Y+\Delta_Y+cF_\lambda,
\]
其中 $F_\lambda\in|F|$。用 $C(\lambda)$ 表示 $\Delta_Y+cF_\lambda$ 中出现的 $f$-例外除子的系数的最大值。因 $|F|$ 为无基点线性系，对一般元 $F_g\in|F|$，$\Delta_Y+cF_g$ 为单纯正规交叉除子。由推论 2.32(2)，
\[
\operatorname{discrep}(X,\Delta+cL_g)=-C(g).
\]
对任意 $F_0\in|F|$，有 $C(0)\ge C(g)$。由定义容易得到不等式 $\operatorname{discrep}(X,\Delta+cL_0)\le-C(0)$。
\end{proof}

% [原文 p.72]
接着，定义极小模型理论中最重要的 6 类奇点中的 5 类。

\begin{definition}{2.34}
设 $(X,\Delta)$ 为正规簇 $X$ 与互不相同素除子之和 $\Delta=\sum a_iD_i$ 的对。（这里系数 $a_i$ 为任意有理数。）假设存在适当的整数 $m>0$，使 $m(K_X+\Delta)$ 为 Cartier 除子。此时按 $\operatorname{discrep}(X,\Delta)$ 的值，称 $(X,\Delta)$ 为终端、标准、klt、plt、lc 中的一种：
\[
\begin{cases}
\operatorname{discrep}(X,\Delta)>0 & \text{终端（terminal）}\\
\operatorname{discrep}(X,\Delta)\ge0 & \text{标准（canonical）}\\
\operatorname{discrep}(X,\Delta)>-1\ \text{且}\ \lfloor\Delta\rfloor\le0 & \text{klt}\\
\operatorname{discrep}(X,\Delta)>-1 & \text{plt}\\
\operatorname{discrep}(X,\Delta)\ge-1 & \text{lc}
\end{cases}
\]
本定义中允许 $a_i<0$，也许有读者感到困惑。这是因为放宽条件并不妨碍讨论一般性质，而且对使用引理 2.30 等的论证极为方便。不过在极小模型纲领中需要 $\Delta\ge0$，第 4 章的分类中也假定 $\Delta\ge0$。
\end{definition}

% [Tutor Note: 这 5 类满足 klt 包含于 plt 包含于 lc，终端包含于标准；越往下允许的奇点越坏，而极小模型纲领要找的正是「坏」与「可计算」之间的那条窄带。]

这里，klt 是「川又对数终端（Kawamata log terminal）」的略语，plt 是「纯对数终端（purely log terminal）」的，lc 是「对数标准（log canonical）」的。（常用的「$(X,\Delta)$ 具有终端奇点」这一说法，也有时是指 $(X,0)$ 的奇点，故容易混淆。）
若 $(X,\Delta=\sum a_iD_i)$ 为 lc（当然 $D_i$ 互不相同），则由推论 2.31(1)，对一切 $i$ 有 $a_i\le1$。
这 5 个概念在极小模型理论中都占有重要位置。
\begin{enumerate}[label=（\arabic*）,leftmargin=2.6em,itemsep=1pt]
\item 终端：若设 $\Delta=0$，这是执行非奇异簇的极小模型纲领所需要的最小的类。$\Delta\ne0$ 的概念只是顺便使用。
\item 标准：若设 $\Delta=0$，这正是具有一般型的簇的标准
% [原文 p.73]
模型上出现的奇点。它对模空间问题特别重要。
\item 川又对数终端：消没定理的证明在这一类中自然地进行。一般说来，在标准奇点上起作用的证明，如第 3 章所见，往往在这一类中也起作用。不过这一类不包含开簇的情形，也不适合归纳证明。
若 $\Delta=0$，则 klt、plt 以及 dlt（定义 2.37）的概念一致；此时称 $X$ 具有对数终端（log terminal，简记 lt）奇点。
\item 纯对数终端：这类是为归纳证明而定义的。本书不常使用。
\item 对数标准：这是使差异系数仍有意义的最大类。遗憾的是，从上同调的观点看，它包含许多复杂情形，故也是不好用的类。第 6 个类是为克服其中若干问题而在定义 2.37 中引入的。
\end{enumerate}
引理 2.27 与推论 2.32 意味着这些概念的某种连续性。

\begin{corollary}{2.35}
设 $(X,\Delta)$ 为对，$\Delta'$ 为有效 $\Q$-Cartier 除子。此时成立。
\begin{enumerate}[label=（\arabic*）,leftmargin=2.6em,itemsep=1pt]
\item 若 $(X,\Delta+\Delta')$ 为终端（或标准、klt、plt、lc），则 $(X,\Delta)$ 也为终端（或标准、klt、plt、lc）。
\item 若 $(X,\Delta)$ 为终端（或 klt），则对 $0\le\varepsilon\ll1$，$(X,\Delta+\varepsilon\Delta')$ 也仍为终端（或 klt）。
\item 若 $(X,\Delta)$ 为 plt，且 $\Delta$ 与 $\Delta'$ 无公共成分，则对 $0\le\varepsilon\ll1$，$(X,\Delta+\varepsilon\Delta')$ 也仍为 plt。
\item 设 $(X,\Delta)$ 为终端。此时 $(X,\Delta+\Delta')$ 为标准，当且仅当对一切 $c<1$，$(X,\Delta+c\Delta')$ 为终端。
\item 若 $(X,\Delta)$ 为 klt（或 plt），则 $(X,\Delta+\Delta')$ 为 lc 当且仅当对一切 $c<1$，$(X,\Delta+c\Delta')$ 为 klt（或 plt）。
\end{enumerate}
\end{corollary}

% [Tutor Note: 推论 2.35 是「开性」的精确表述：把 $\Delta$ 略微扰动，终端性与 klt 性保持不变。这类开性在归纳证明中不可缺，它保证允许的奇点类不会因微小形变而跳出去。]

% [原文 p.74]
\begin{propositionp}{2.36}
设 $(X,\Delta)$ 为 klt 对。
\begin{enumerate}[label=（\arabic*）,leftmargin=2.6em,itemsep=1pt]
\item $(X,\Delta)$ 的对数奇点消解 $f\colon X'\to X$ 中，存在这样的一个：把 $f^*(K_X+\Delta)\equiv K_{X'}+A_{X'}-B_{X'}$（$A_{X'}$、$B_{X'}$ 为无公共成分的有效除子）表出时，$\Supp A_{X'}$ 非奇异。
\item 若 $a(E,X,\Delta)<1+\operatorname{totaldiscrep}(X,\Delta)$，则 $\operatorname{center}_{X'}E$ 为除子。因此 $X$ 上空的例外除子 $E$ 中满足 $a(E,X,\Delta)<1+\operatorname{totaldiscrep}(X,\Delta)$ 的至多只有有限个。
\end{enumerate}
\end{propositionp}
\begin{proof}
为证 (1)，取任意对数奇点消解 $f\colon Y\to X$，把 $A_Y$ 分解到不同的素除子：$A_Y=\sum_{i\in I}(1-\alpha_i)A_i$（$0<\alpha_i<1$）。令 $S(Y)=\sum_{i\in I}\Z_{\ge0}\alpha_i\subset\R_{\ge0}$。再如下定义。先令
\[
w(J)=\sum_{i\in J}\alpha_i\quad(J\subset I),
\]
当 $\bigcap_{i\in J}A_i\ne\varnothing$ 且 $\#J\ge2$ 时以此式取值；否则令 $w(J)=\infty$。又令
\[
M(Y)=\min\{w(J):J\subset I\},\qquad r(Y)=\#\{J\subset I:w(J)=M(Y)\}.
\]
这里注意以下几点：$M(Y)\in(0,\dim X)\cap S(Y)$，或 $M(Y)=\infty$；并且 $M(Y)=\infty$ 当且仅当 $\Supp A_Y$ 非奇异。
假设 $M(Y)<\infty$。选取 $J\subset I$ 使 $w(J)=M(Y)$，令 $Z=\bigcap_{i\in J}A_i$，$p\colon Y'\to Y$ 为沿 $Z$ 的爆破，$E$ 为例外除子。则 $(f\circ p)^*(K_X+\Delta)$ 中 $E$ 的系数等于 $1-\sum_{i\in J}\alpha_i$，因此 $S(Y')=S(Y)$。进一步容易看出：或者 $M(Y')>M(Y)$，或者 $M(Y')=M(Y)$ 且 $r(Y')<r(Y)$。因 $M(Y)$ 只取有限多个值，爆破不可能永远继续下去。故经有限次爆破后 $M(Y)=\infty$ 成立。
由推论 2.31，若 $E$ 为 $X'$ 上空的例外除子，则
\[
a(E,X,\Delta)\ge a(E,X',A_{X'})\ge\min\bigl\{1,\ \min_i\{\alpha_i\}\bigr\}
\]
% [原文 p.75]
成立；其中 $A_{X'}=\sum_{i\in I}(1-\alpha_i)A_i$。这里由定义与引理 2.30，
\[
\min\bigl\{1,\min_i\{\alpha_i\}\bigr\}\ge1+\operatorname{totaldiscrep}(Y,A_Y-B_Y)=1+\operatorname{totaldiscrep}(X,\Delta).
\]
因此由 (1) 得到 (2)。
\end{proof}

\begin{definition}{2.37}
设 $(X,\Delta)$ 为正规簇 $X$ 与互不相同素除子之和 $\Delta=\sum a_iD_i$ 的对。不过系数 $a_i$ 满足 $0\le a_i\le1$。对适当的整数 $m>0$，$m(K_X+\Delta)$ 为 Cartier 除子。若存在满足下列条件的闭集合 $Z\subset X$，则称对 $(X,\Delta)$ 为 dlt，或除子对数终端（divisional log terminal）。
\begin{enumerate}[label=（\arabic*）,leftmargin=2.6em,itemsep=1pt]
\item $X\setminus Z$ 非奇异，且 $\Delta|_{(X\setminus Z)}$ 为单纯正规交叉除子。
\item 若 $f\colon Y\to X$ 为双有理态射，$E\subset Y$ 为满足 $\operatorname{center}_XE\subset Z$ 的素除子，则 $a(E,X,\Delta)>-1$。
\end{enumerate}
\end{definition}

\begin{remark}{2.38}
这个定义极为微妙。(1) 中必须要求 $\Delta|_{(X\setminus Z)}$ 为「单纯」正规交叉除子。若允许 $\Delta$ 有自交，这个更一般的类也许会表现良好，但许多证明会行不通。
定义 2.37 与通常的定义（参见 [K+92, 2.13]）不同，但它们由定理 2.44 等价。之所以这样定义，是因为定理 2.44 要用 [Hir64] 的手法巧妙地证明，而从本书的定义出发，不用定理 2.44 也能证明所需的一切结果。
另外，本书（§3.7）主要讨论 dlt 对的极小模型纲领。因为极小模型纲领需要 $\Delta\ge0$，故定义 2.37 也假定 $\Delta\ge0$。
\end{remark}

% [Tutor Note: 注意 2.38 抓住的正是 dlt 与一般 lc 的分水岭：允许自交就丢掉了「单纯」正规交叉，而消没定理与归纳法都依赖这一条。]

接下来叙述的 dlt 对的连续性对应于推论 2.35，由引理 2.27 与推论 2.32 同样得到。

\begin{corollary}{2.39}
设 $(X,\Delta)$ 为 $\Delta$ 有效的对，$\Delta'$ 为有效 $\Q$-Cartier 除子。此时：
\begin{enumerate}[label=（\arabic*）,leftmargin=2.6em,itemsep=1pt]
\item 若 $(X,\Delta+\Delta')$ 为 dlt，则 $(X,\Delta)$ 也仍为 dlt。
\item 若 $(X,\Delta)$ 为 dlt，且同时 $\Supp\Delta'\subset\Supp(\Delta-\lfloor\Delta\rfloor)$，则对 $0\le\varepsilon\ll1$，$(X,\Delta+\varepsilon\Delta')$ 也仍为 dlt。
\end{enumerate}
\end{corollary}

% [原文 p.76]
\begin{propositionp}{2.40}
考虑正规簇 $X$ 与 $\Q$-除子 $\Delta$ 的对 $(X,\Delta)$。不过 $\Delta=\sum a_iD_i$（互不相同素除子之和）的系数满足 $0\le a_i\le1$。设对适当的整数 $m>0$，$m(K_X+\Delta)$ 为 Cartier 除子。此时，$(X,\Delta)$ 为 dlt 的充要条件是：存在满足下列 2 个条件的闭集合 $Z\subset X$。
\begin{enumerate}[label=（\arabic*）,leftmargin=2.6em,itemsep=1pt]
\item $X\setminus Z$ 非奇异，且 $\Delta|_{(X\setminus Z)}$ 为单纯正规交叉除子。
\item $(X,\Delta)$ 的对数奇点消解 $f\colon Y\to X$ 中，存在这样的一个：$f^{-1}(Z)\subset Y$ 为纯余维数 $1$，且对一切不可约除子 $E\subset f^{-1}(Z)$ 有 $a(E,X,\Delta)>-1$。
\end{enumerate}
\end{propositionp}
\begin{proof}
若 $(X,\Delta)$ 为 dlt，则满足 $f^{-1}(Z)\subset Y$ 为纯余维数 $1$ 的任意对数奇点消解 $f$ 都具有上述性质。
反之，设存在具有上述性质的 $f\colon Y\to X$。用 $K_Y+\Delta_Y\equiv f^*(K_X+\Delta)$ 表出，并设 $\Delta'$ 为满足 $\Supp\Delta'=\Supp f^{-1}(Z)$ 的有效除子。由性质 (2)，$\Delta'$ 的任意不可约成分在 $\Delta_Y$ 中的系数 $<1$。因此对 $0<\varepsilon\ll1$，$\Delta_Y+\varepsilon\Delta'$ 的任意不可约成分的系数 $\le1$。故由推论 2.31(3)，$(Y,\Delta_Y+\varepsilon\Delta')$ 为 lc。
设 $E$ 为中心含于 $Z$ 的 $X$ 上空任意除子。此时 $\operatorname{center}_YE\subset f^{-1}(Z)$，故由引理 2.27，
\[
a(E,X,\Delta)=a(E,Y,\Delta_Y)>a(E,Y,\Delta_Y+\varepsilon\Delta')\ge-1.
\]
这样定义 2.37(2) 的条件就满足了。
\end{proof}

\begin{propositionp}{2.41}
设 $(X,\Delta)$ 为 dlt 对。此时它为 klt 的充要条件是 $\lfloor\Delta\rfloor=0$。
\end{propositionp}
\begin{proof}
若 $(X,\Delta)$ 为 klt，则在 dlt 的定义 2.37 中取 $Z=X$，即知它为 dlt。
设 $(X,\Delta)$ 为 dlt 且 $\lfloor\Delta\rfloor=0$。设 $f\colon Y\to X$ 为双有理态射，$E\subset Y$ 为例外除子。若 $\operatorname{center}_XE\subset Z$，则由 dlt 的定义有 $a(E,X,\Delta)>-1$。若 $\operatorname{center}_XE\not\subset Z$，则由推论 2.31(3) 有 $a(E,X,\Delta)>-1$。于是 $(X,\Delta)$ 为 klt。
\end{proof}

\begin{propositionp}{2.42}
设 $(X,\Delta)$ 为曲面上的 dlt 对，$x\in X$。此时在 $x$ 的适当邻域 $U$ 上，或者 $X$ 非奇异、$\Delta=\lfloor\Delta\rfloor$ 具有横断相交的 2 个不可约成分，或者 $(X,\Delta)$ 为 plt；两者必居其一。
\end{propositionp}
